% TOP_PROBLEM: {"id": 30007184, "title": "Nonexpansive fixed points in reflexive Banach spaces", "category_id": 8, "category": {"id": 8, "name": "analysis", "display_name": "Analysis", "description": "Limits, continuity, calculus, and function theory.", "slug": "analysis", "order_index": 8, "created_at": "2026-07-31T15:26:25.671Z"}, "status": "open"}
% FIRST_POSTED: 2026-10-08
% ATTEMPT_STATUS: unresolved
\documentclass[11pt]{article}
\usepackage[a4paper,margin=27mm]{geometry}
\usepackage[T1]{fontenc}
\usepackage[utf8]{inputenc}
\usepackage{lmodern}
\usepackage{amsmath,amssymb,amsthm,mathtools,enumitem,booktabs,array,microtype}
\usepackage[colorlinks=true,urlcolor=blue,citecolor=blue,linkcolor=blue]{hyperref}
\usepackage{cleveref}
\setlength{\emergencystretch}{3em}
\Urlmuskip=0mu plus 1mu\relax
\newtheorem{theorem}{Theorem}[section]
\newtheorem{proposition}[theorem]{Proposition}
\newtheorem{corollary}[theorem]{Corollary}
\newtheorem{lemma}[theorem]{Lemma}
\newtheorem{remark}[theorem]{Remark}
\theoremstyle{definition}\newtheorem{definition}[theorem]{Definition}
\newcommand{\Fix}{\operatorname{Fix}}
\newcommand{\BR}{\operatorname{BR}}
\newcommand{\argmax}{\operatorname*{arg\,max}}
\newcommand{\argmin}{\operatorname*{arg\,min}}
\newcommand{\R}{\mathbb R}
\newcommand{\N}{\mathbb N}
\title{Nonexpansive Fixed Points and Equilibrium Existence in Reflexive Banach Spaces\\\large Nash Equilibria, Strategic Networks, Aggregative Games, and Equilibrium-Set Geometry}


\begin{document}

\section{Foundational hypothesis, provenance, and boundaries}
\begin{theorem}[Reflexive nonexpansive fixed-point assertion (FP)\label{thm:FP}]
Let $X$ be a real reflexive Banach space with its given norm and let $C\subset X$ be nonempty, bounded, norm-closed and convex. If $T:C\to C$ satisfies
\begin{equation}\label{eq:nonexp}\|Tx-Ty\|\le\|x-y\|\quad(x,y\in C),\end{equation}
then $\Fix(T)\ne\varnothing$.
\end{theorem}
This is the main theorem asserted by the OpenAI manuscript dated 24 September 2026 \cite{OAIFP}, identified as problem LOCAL-30007184 by the tracking website \cite{Tracker}. The released repository supplies an implementation of its main fixed-point theorem in \texttt{lean/OAI/Analysis/Nonexpansive/Main.lean}, together with a comparator challenge and scope note \cite{LeanMain,LeanScope,LeanComparator}. The comparator challenge's \texttt{sorry} specifies the verification target, rather than being the implementation of that main declaration. Repository inspection is \emph{not} a substitute for rerunning Lean or checking every imported dependency. All the results in this note are proved \emph{conditional on \cref{thm:FP}}; claims about the independent verification of \cref{thm:FP} are intentionally not made.

Classical compact-convex equilibrium theorems \cite{Glicksberg1952}, fixed points for uniformly convex spaces \cite{Browder1965,Gohde1965}, and reflexive spaces with normal structure \cite{Kirk1965} already imply numerous special cases. The distinguishing stated generality of \cref{thm:FP} is that neither norm compactness, weak continuity of $T$, nor an additional geometric property of the given reflexive norm is imposed. Norm continuity does \emph{not} in general imply weak-to-weak continuity. Thus reflexivity alone cannot justify invoking a weakly continuous fixed-point theorem for an arbitrary norm-nonexpansive map. None of the following results assert equilibrium for arbitrary games with continuous, concave payoffs.

\begin{definition}[Selected best-response map]
For a game with players $I$, feasible profiles $C=\prod_{i\in I}C_i$, and payoffs $u_i$, a \emph{selected best-response map} is a function $B:C\to C$ such that, for every $x\in C$ and every $i$,
\[B_i(x)\in \argmax_{a_i\in C_i}u_i(a_i,x_{-i}).\]
Its fixed points are Nash equilibria. When each best response is unique, its fixed-point set equals the entire Nash equilibrium set; with nonunique responses, it may be a proper subset.
\end{definition}

\section{General Nash existence and exact logical strength}
\begin{theorem}[Finite-player Nash existence\label{thm:finite}]
Let $I=\{1,\ldots,n\}$, let each $X_i$ be a real reflexive Banach space, and let each $C_i\subset X_i$ be nonempty, bounded, norm-closed and convex. Suppose the game admits a selected best-response map $B:C=\prod_i C_i\to C$ that is nonexpansive for some norm $\|\cdot\|_*$ equivalent to a standard finite-product norm. Then the game admits a pure Nash equilibrium.
\end{theorem}
\begin{proof}
The product $X=\prod_iX_i$, equipped with $\|\cdot\|_*$, is a reflexive Banach space because a finite product of reflexive spaces is reflexive and an equivalent renorming preserves reflexivity. Its subset $C$ is nonempty, bounded, closed and convex. By \cref{thm:FP}, $B(x^*)=x^*$ for some $x^*\in C$. Each coordinate satisfies $x_i^*\in\argmax_{a_i\in C_i}u_i(a_i,x_{-i}^*)$, the Nash inequalities.
\end{proof}
\begin{remark}
A selected response is assumed to \emph{exist}; neither nonemptiness of best-response sets nor their convexity guarantees a nonexpansive selection. Moreover nonexpansiveness must hold for a specified joint norm, not merely separate informal Lipschitz bounds. No independent existence claim is made for general correspondences.
\end{remark}

\begin{theorem}[Exact equivalence with two-player tracking games\label{thm:equivalence}]
For a real Banach space $X$, the following are equivalent:
\begin{enumerate}[label=(\roman*)]
\item Every nonexpansive self-map of every nonempty bounded closed convex $D\subset X$ has a fixed point.
\item For each such $D$ and $T:D\to D$, the two-player game with both action sets $D$ and payoffs
\[u_1(x,y)=-\|x-Ty\|^2,\qquad u_2(x,y)=-\|y-x\|^2\]
has a pure Nash equilibrium.
\end{enumerate}
If $X$ is reflexive, both properties follow from \cref{thm:FP}.
\end{theorem}
\begin{proof}
For a fixed $y$, the unique maximizer of $u_1(\cdot,y)$ is $Ty\in D$. For fixed $x$, the unique maximizer of $u_2(x,\cdot)$ is $x\in D$. Therefore $(x,y)$ is Nash precisely when $x=Ty$ and $y=x$, which is equivalent to $x=y=Tx$. This is a two-way equivalence for every $D,T$. Both payoffs are norm-continuous and concave in own action because the square of a norm is convex. Under reflexivity the fixed-point side holds by \cref{thm:FP}.
\end{proof}
\begin{remark}[Interpretation of novelty]
\Cref{thm:equivalence} imports the full logical generality of the fixed-point claim into a concrete class of games, but does not demonstrate that a specific previously named economic open problem is solved. Many special cases are already classical. It is also important that the tracking construction does not require differentiable norms or Hilbert structure.
\end{remark}

\section{Strategic interaction networks and primitive conditions}
\begin{theorem}[Finite network spectral-radius criterion\label{thm:spectral}]
Retain the strategy-domain assumptions of \cref{thm:finite}. Suppose every player has a unique best-response function $b_i(x_{-i})$ and there is a nonnegative $n\times n$ matrix $A=(a_{ij})$ with $a_{ii}=0$ such that
\begin{equation}\label{eq:matrixlip}
\|b_i(x_{-i})-b_i(y_{-i})\|_i\le\sum_{j=1}^n a_{ij}\|x_j-y_j\|_j
\quad(i=1,\ldots,n).
\end{equation}
If $\rho(A)\le1$, the game has a Nash equilibrium. If $\rho(A)<1$, its Nash equilibrium is unique.
\end{theorem}
\begin{proof}
First suppose $A$ is irreducible. By Perron--Frobenius, some $w\gg0$ satisfies $Aw=\rho(A)w$. In the weighted maximum norm $\|x\|_w=\max_i\|x_i\|_i/w_i$, inequality \eqref{eq:matrixlip} yields
\[\|B(x)-B(y)\|_w\le\rho(A)\|x-y\|_w.\]
Existence follows from \cref{thm:finite} when $\rho(A)\le1$.

For reducible $A$, topologically sort the strongly connected components of its directed dependence graph. The resulting matrix is block triangular, each diagonal block having spectral radius at most $\rho(A)\le1$. For a nontrivial irreducible block, its own response map (holding previously solved blocks fixed) is nonexpansive in its Perron-weighted maximum norm. For a singleton block, $a_{ii}=0$, so the corresponding response is constant once earlier dependencies are fixed. Solve the blocks successively by \cref{thm:FP}; later components cannot affect earlier responses because of the block-triangular dependence. This constructs a Nash profile.

If $\rho(A)<1$, let $w=(I-A)^{-1}\mathbf1=\sum_{k\ge0}A^k\mathbf1\gg0$. Then $Aw=w-\mathbf1<w$, and hence $q=\max_i(Aw)_i/w_i<1$. The global joint response is a $q$-contraction in $\|\cdot\|_w$. Banach's contraction principle implies its fixed point is unique, and the uniqueness of individual best responses identifies it with the game's unique Nash equilibrium.
\end{proof}
\begin{remark}
The borderline $\rho(A)=1$ permits no strict contraction. A \emph{single} global positive super-eigenvector $Aw\le w$ need not exist for reducible $A$ at the boundary; the block argument above is essential. For two players, the criterion is $a_{12}a_{21}\le1$.
\end{remark}

\begin{theorem}[Strong own-action convexity versus cross-effects\label{thm:primitives}]
Let each player $i$ minimize $c_i(a_i,x_{-i})$ over $C_i$, with the spaces and domains of \cref{thm:finite}. Assume that for every $x_{-i}$, the map $a_i\mapsto c_i(a_i,x_{-i})$ is proper, finite-valued, norm-lower-semicontinuous, $m_i$-strongly convex ($m_i>0$), and Fr\'echet differentiable on a neighborhood of $C_i$. Assume for every $a_i\in C_i$ and all opponents' profiles that
\[\|D_i c_i(a_i,x_{-i})-D_i c_i(a_i,y_{-i})\|_{X_i^*}
\le\sum_{j\ne i}L_{ij}\|x_j-y_j\|_j,\quad L_{ij}\ge0.\]
Put $A_{ii}=0$ and $A_{ij}=L_{ij}/m_i$ for $j\ne i$. If $\rho(A)\le1$, a pure Nash equilibrium exists; if $\rho(A)<1$, it is unique.
\end{theorem}
\begin{proof}
Each $C_i$ is weakly compact. Convex norm-lower-semicontinuous costs are weakly lower semicontinuous, so a minimizer exists; strong convexity makes it unique. Denote the minimizers against $x_{-i}$ and $y_{-i}$ by $a$ and $b$. First-order optimality gives
\[\langle D_ic_i(a,x_{-i}),b-a\rangle\ge0,\quad
\langle D_ic_i(b,y_{-i}),a-b\rangle\ge0.\]
Strong monotonicity and the two variational inequalities imply
\begin{align*}
m_i\|a-b\|_i^2
&\le\langle D_ic_i(a,x_{-i})-D_ic_i(b,x_{-i}),a-b\rangle\\
&\le\langle D_ic_i(b,y_{-i})-D_ic_i(b,x_{-i}),a-b\rangle\\
&\le\sum_{j\ne i}L_{ij}\|x_j-y_j\|_j\,\|a-b\|_i.
\end{align*}
Cancel $\|a-b\|_i$ if it is nonzero to obtain \eqref{eq:matrixlip} for $a_{ij}=L_{ij}/m_i$. Apply \cref{thm:spectral}.
\end{proof}
\begin{remark}
We use the convention that $m$-strong convexity implies $\langle Dc(a)-Dc(b),a-b\rangle\ge m\|a-b\|^2$. The condition compares marginal-cost curvature to cross-player sensitivity; it requires neither a potential game nor strategic complementarity.
\end{remark}

\section{Games with infinitely many players}
\begin{theorem}[Countably many players in a reflexive direct sum\label{thm:countable}]
Let $1<p<\infty$ and let $(X_i)_{i\ge1}$ be real reflexive Banach spaces. Each $C_i\subset X_i$ is nonempty, closed and convex, and $\sup_{a\in C_i}\|a\|_i\le r_i$. Choose $\eta_i>0$ with $\sum_i\eta_i r_i^p<\infty$. On the weighted direct sum
\[X=\Bigl(\bigoplus_{i\ge1}X_i\Bigr)_{\ell^p(\eta)},\qquad
\|x\|_{p,\eta}=\Bigl(\sum_i\eta_i\|x_i\|_i^p\Bigr)^{1/p},\]
let $C=\prod_iC_i\subset X$. Suppose a selected best-response map $B:C\to C$ obeys
\[\|B_i(x)-B_i(y)\|_i\le\sum_j a_{ij}\|x_j-y_j\|_j\]
where $A=(a_{ij})$ is nonnegative and defines a bounded linear positive operator on $\ell^p(\eta)$ of norm at most one. Then the countable-player game has a pure Nash equilibrium.
\end{theorem}
\begin{proof}
Weighted $\ell^p$-sums of reflexive spaces are reflexive. The uniform summability condition places the entire product $C$ inside a bounded subset of $X$. Coordinate projections are continuous, so $C=\bigcap_i\pi_i^{-1}(C_i)$ is norm-closed; it is also nonempty and convex. Writing $d_j=\|x_j-y_j\|_j$, the assumed inequalities imply
\[\|B(x)-B(y)\|_{p,\eta}\le\|Ad\|_{\ell^p(\eta)}\le\|d\|_{\ell^p(\eta)}=\|x-y\|_{p,\eta}.\]
Apply \cref{thm:FP} to $B:C\to C$. At its fixed point every coordinate is a best response.
\end{proof}
\begin{remark}
The whole Cartesian product is admissible here because of the summable radius envelope. An arbitrary unbounded or $\ell^\infty$-product is not covered. In certain finite-dimensional-coordinate cases, these assumptions already make $C$ norm-compact, so the existence conclusion is then classical.
\end{remark}

\begin{theorem}[Nonatomic aggregative Nash equilibrium\label{thm:nonatomic}]
Let $(S,\Sigma,\mu)$ be an atomless probability space, $E$ a separable real reflexive Banach space, $1<p<\infty$, and $Z$ a Banach space. For each $s\in S$, let $K(s)\subset E$ be nonempty, closed and convex, with $\sup_{a\in K(s)}\|a\|\le h(s)$ almost everywhere, where $h\in L^p(\mu)$. Suppose the feasible set
\[\mathcal C=\{f\in L^p(S;E):f(s)\in K(s)\text{ a.e.}\}\]
is nonempty. Let $\mathcal A:L^p(S;E)\to Z$ be bounded linear. Assume measurable individual best responses $r(s,z)\in\argmax_{a\in K(s)}v(s,a,z)$ exist and satisfy, outside one common null set,
\[\|r(s,z)-r(s,z')\|\le\kappa(s)\|z-z'\|_Z,\qquad \kappa\in L^p(\mu).\]
If $\|\kappa\|_{L^p}\|\mathcal A\|\le1$, there exists $f^*\in\mathcal C$ with
\[f^*(s)\in\argmax_{a\in K(s)}v(s,a,\mathcal Af^*)\quad\text{for a.e. }s.\]
\end{theorem}
\begin{proof}
The space $L^p(S;E)$ is reflexive. The envelope $h$ bounds $\mathcal C$ and gives integrability of $r(\cdot,z)$, while convexity is pointwise. To see that $\mathcal C$ is norm-closed, take an $L^p$-convergent sequence in it, extract an almost-everywhere pointwise convergent subsequence, and invoke closedness of $K(s)$. Set $(Bf)(s)=r(s,\mathcal Af)$. Strong measurability and the envelope make this a well-defined element of $\mathcal C$. By the pointwise Lipschitz bound and linearity of $\mathcal A$,
\[\|Bf-Bg\|_{L^p}\le\|\kappa\|_{L^p}\|\mathcal Af-\mathcal Ag\|_Z
\le\|\kappa\|_{L^p}\|\mathcal A\|\|f-g\|_{L^p}\le\|f-g\|_{L^p}.\]
A fixed point from \cref{thm:FP} yields the desired almost-everywhere optimality.
\end{proof}
\begin{remark}
This is a nonatomic game in which each player takes the aggregate as given. The theorem does not license treating the aggregate as exogenous in a finite-player model where unilateral decisions materially change it. The choice $\mathcal Af=\int_Sf\,d\mu$ is a canonical example. Measurability and the existence of optimal response rules are essential assumptions.
\end{remark}

\section{Generalized Nash equilibrium and aggregate consistency}
\begin{theorem}[Interdependent constraints\label{thm:generalized}]
Retain the finite reflexive domains $C_i$ from \cref{thm:finite}. Suppose player $i$ faces a nonempty feasible set $K_i(x_{-i})\subseteq C_i$ conditional on opponents' actions. If a jointly nonexpansive map $B:C\to C$ exists in an equivalent product norm and satisfies
\[B_i(x)\in K_i(x_{-i}),\qquad
u_i(B_i(x),x_{-i})\ge u_i(a_i,x_{-i})\quad
( a_i\in K_i(x_{-i})),\]
where $u_i$ denotes player $i$'s utility, then a generalized Nash equilibrium exists.
\end{theorem}
\begin{proof}
At a fixed point of $B$ supplied by \cref{thm:FP}, every coordinate is feasible and optimal against the other coordinates. These are exactly the generalized Nash conditions.
\end{proof}
\begin{remark}
The assumption requires a globally defined, nonexpansive \emph{feasible optimal-response selection}. It does not follow just from convexity of the constraint correspondences.
\end{remark}

\begin{corollary}[Aggregate response consistency\label{cor:aggregate}]
Let $M$ be a nonempty bounded closed convex subset of a real reflexive Banach space $Z$. Suppose $R:M\to Y$ specifies individual optimal behavior and $\mathcal A:R(M)\to M$ aggregates these choices, with $T=\mathcal A\circ R:M\to M$ nonexpansive. Then there is $m^*\in M$ such that $m^*=\mathcal A(R(m^*))$.
\end{corollary}
\begin{proof}
Apply \cref{thm:FP} to $T$.\end{proof}
\begin{remark}
The conclusion is only an aggregate consistency condition. It becomes an economic equilibrium if, in addition, the constructions of $R$ and $\mathcal A$ incorporate the model's optimization, information, feasibility and any required market-clearing conditions. Neither factor needs individually to have Lipschitz constant at most one as long as the composite is nonexpansive.
\end{remark}

\section{Equilibrium-set retractions, selection, and common equilibria}
The source manuscript further asserts a corollary derived from Bruck's theorem \cite{OAIFP,Bruck1974}: for a nonempty bounded closed convex subset $C$ of a reflexive Banach space, the common fixed-point set of any pairwise commuting family of nonexpansive self-maps is nonempty and a \emph{nonexpansive retract} of $C$. We use that corollary as a second explicit source input, denoted (BR). In particular, it is not derived merely from the bare existence statement (FP) without an additional fixed-point-set theorem.
\begin{theorem}[Nonexpansive equilibrium selection and contractibility\label{thm:retraction}]
Under the conditions of \cref{thm:finite}, assume additionally the commuting-family retraction corollary (BR) stated above. Set $E_B=\Fix(B)$. There exists a nonexpansive retraction $P:C\to E_B$ such that $P(e)=e$ for all $e\in E_B$. Consequently:
\begin{enumerate}[label=(\alph*)]
\item for every reference profile $a\in C$, $P(a)$ is a selected equilibrium and $\|P(a)-P(a')\|_*\le\|a-a'\|_*$;
\item $E_B$ is nonempty and contractible in the norm topology, hence path-connected.
\end{enumerate}
If every player's best response is unique, $E_B$ is the full Nash equilibrium set.
\end{theorem}
\begin{proof}
Apply (BR) to the singleton commuting family $\{B\}$. For $e_0,e_1\in E_B$, the map $t\mapsto P((1-t)e_0+te_1)$ is a continuous path through $E_B$ joining the endpoints. Fix $e_0\in E_B$ and define
\[H:E_B\times[0,1]\to E_B,\qquad H(e,t)=P((1-t)e+te_0).\]
Then $H(e,0)=e$, $H(e,1)=e_0$, and $H(e_0,t)=e_0$. Continuity follows from continuity of affine interpolation and the Lipschitz property of $P$. Thus $E_B$ is contractible. The selection bound follows directly from nonexpansiveness of $P$.
\end{proof}
\begin{corollary}[Compatible scenarios\label{cor:commuting}]
Suppose maps $B_\theta:C\to C$ are each nonexpansive, are pairwise commuting, and represent selected optimal-response rules for a collection of games or scenarios with the same feasible profile set $C$. Under (BR) there is $x^*\in C$ satisfying $B_\theta(x^*)=x^*$ for every $\theta$; their common selected-equilibrium set is a nonempty nonexpansive retract of $C$.
\end{corollary}
\begin{proof}
Apply (BR) to the commuting family $\{B_\theta\}_\theta$ and translate each fixed-point equation into the relevant Nash inequalities.
\end{proof}
\begin{remark}
The retraction need not be a metric projection, economically canonical or computationally tractable. Its stability compares \emph{reference profiles within one game}, not necessarily equilibria across changes in economic primitives. Pairwise commuting responses across scenarios are a strong assumption; unilateral best-response updates generally need not commute.
\end{remark}

\section{Sharpness, dynamics, and scope limitations}
\begin{proposition}[No Lipschitz relaxation beyond one\label{prop:sharp}]
For every $\varepsilon>0$ there exists a two-player game on the closed unit ball $D$ of the Hilbert space $\ell^2$, with jointly norm-continuous payoffs that are strongly concave in own actions, such that its selected best-response map is Lipschitz with constant strictly less than $1+\varepsilon$ under a natural equivalent product norm, but the game has no Nash equilibrium.
\end{proposition}
\begin{proof}
Fix $0<\delta\le1$ such that $\sqrt{1+\delta^2}<1+\varepsilon$ and define
\[T_\delta(x_1,x_2,\ldots)=\bigl(\delta(1-\|x\|),x_1,x_2,\ldots\bigr),\quad x\in D.\]
For $r=\|x\|\in[0,1]$, $\|T_\delta x\|^2=\delta^2(1-r)^2+r^2\le1$, so $T_\delta(D)\subset D$. For all $x,y$,
\[\|T_\delta x-T_\delta y\|^2=\delta^2(\|x\|-\|y\|)^2+\|x-y\|^2
\le(1+\delta^2)\|x-y\|^2.\]
A fixed point would satisfy $x_{k+1}=x_k$ for every $k$, so square-summability forces $x=0$, contradicting the first-coordinate equation $x_1=\delta$. Define
\[u_1(x,y)=-\|x-T_\delta y\|^2,\qquad u_2(x,y)=-\|y-x\|^2.\]
Unique best responses are $T_\delta y$ and $x$, so an equilibrium would give a fixed point of $T_\delta$, which is impossible. In the maximum product norm, the joint response $B(x,y)=(T_\delta y,x)$ has Lipschitz constant at most $\sqrt{1+\delta^2}<1+\varepsilon$. Its payoffs are continuous and strongly concave in the respective own choices in the Hilbert norm.
\end{proof}

\begin{proposition}[Existence does not ensure best-response convergence\label{prop:cycles}]
There is a game with a unique Nash equilibrium and a nonexpansive joint response map for which every nonzero simultaneous best-response orbit is periodic of period four.
\end{proposition}
\begin{proof}
Let $D$ be the unit ball of any real reflexive Banach space and let both players have action set $D$. Choose $u_1(x,y)=-\|x+y\|^2$ and $u_2(x,y)=-\|y-x\|^2$. Their unique best responses yield $B(x,y)=(-y,x)$, an isometry for the maximum product norm. The unique fixed point is $(0,0)$, whereas $B^4=\mathrm{id}$ and $B^2=-\mathrm{id}$. Every nonzero profile consequently cycles with minimal period four.
\end{proof}

\begin{proposition}[Approximate equilibria from anchored contractions\label{prop:approx}]
Let $C$ and nonexpansive $B:C\to C$ satisfy the hypotheses of \cref{thm:FP}. Fix $a\in C$. For every $0<\epsilon<1$, the map $T_\epsilon=(1-\epsilon)B+\epsilon a$ has a unique fixed point $x_\epsilon\in C$, satisfying
\[\|x_\epsilon-Bx_\epsilon\|\le\epsilon\operatorname{diam}(C).\]
If, in a finite-player game, each $u_i(\cdot,x_{-i})$ is $L_i$-Lipschitz on $C_i$ uniformly in $x_{-i}$, and the product norm dominates the $i$th coordinate with constant $d_i$, then $x_\epsilon$ is a unilateral-regret approximate Nash equilibrium with regret for player $i$ at most $L_i d_i\epsilon\operatorname{diam}(C)$.
\end{proposition}
\begin{proof}
$T_\epsilon$ is an $(1-\epsilon)$-contraction on the complete set $C$, so the contraction theorem supplies its unique fixed point. Its equation gives $x_\epsilon-Bx_\epsilon=\epsilon(a-Bx_\epsilon)$; both $a$ and $Bx_\epsilon$ lie in $C$. For player $i$, let $b_i=B_i(x_\epsilon)$ be an optimal response. For any $z_i\in C_i$,
\[u_i(z_i,x_{\epsilon,-i})-u_i(x_{\epsilon,i},x_{\epsilon,-i})\le u_i(b_i,x_{\epsilon,-i})-u_i(x_{\epsilon,i},x_{\epsilon,-i})\le L_i\|b_i-x_{\epsilon,i}\|_i,\]
which is bounded by the displayed coordinate estimate.
\end{proof}
\begin{remark}
The fixed-point theorem by itself does not imply convergence of $B^k(x)$, nor that weak cluster points of approximate fixed points are fixed points: norm-continuous maps may fail to be weakly continuous. It does not address arbitrary correspondences; general probability-measure strategy sets; nonreflexive spaces such as standard $L^1$, $L^\infty$, or $C([0,1])$; welfare efficiency; or numerical computation of a fixed point. A mixed-strategy model restricted to an appropriate reflexive class of densities is a restricted game unless deviations in the full strategy class are also controlled.
\end{remark}


\section{Economic interpretation and implications}
\subsection{Economic intuition, mechanisms, and interpretative limits}\label{subsec:econintuition}

\paragraph{Mutual optimality as a fixed point.}
The mathematical object in the foregoing theorems is the \emph{joint optimal-response operator} $B$. An equilibrium $x^*$ is a profile that reproduces itself when every participant optimizes given the actions of the others: $B(x^*)=x^*$. This is the familiar strategic consistency underlying Nash equilibrium, not a new equilibrium concept. What changes under the reflexive fixed-point assertion (FP) is the collection of response operators for which such self-consistency is guaranteed. The crucial distinction is between a \emph{nonexpansive} response, which does not magnify distances, and a \emph{contractive} response, which reduces them by a uniform factor strictly below one. Under nonexpansiveness, a consistent profile exists even at the knife-edge at which iterative responses need not approach it. A selected best-response map remains an assumption of the model: (FP) does not ensure that maximizers exist or that a multivalued correspondence admits a nonexpansive selection.

\paragraph{Strategic complements, substitutes, and feedback loops.}
The coefficients $a_{ij}$ in \cref{thm:spectral} are upper bounds on the \emph{magnitude} of player $i$'s optimal adjustment when player $j$ changes an action. They therefore accommodate both strategic complementarity and substitutability; their nonnegativity does not mean that individual actions are strategic complements. The spectral radius $\rho(A)$ captures amplification around chains and cycles of interdependence. In the two-player case, $a_{12}a_{21}\le1$ permits one response to amplify deviations so long as the other response compensates. At $\rho(A)<1$, a weighted norm makes joint responses contractive, providing both equilibrium uniqueness and convergence of simultaneous best-response iteration. At $\rho(A)=1$, existence may persist but neither property follows. The reducible-network extension is economically relevant where firms, sectors, or jurisdictions influence one another in a directional hierarchy: strongly connected groups can be solved successively, even where a single global positive system of nonexpansiveness weights is unavailable.

\paragraph{Own incentives versus strategic spillovers.}
The primitive criterion of \cref{thm:primitives} expresses the same economic force without assuming a Lipschitz estimate for responses outright. Strong convexity of cost in one's own choice (or strong concavity of utility) is a restoring force: it penalizes moving far from one's optimum. The cross-effect constant $L_{ij}$ measures how strongly another player's behavior shifts that player's marginal incentives. The ratio $L_{ij}/m_i$ expresses cross-player sensitivity relative to own-action curvature. The network criterion $\rho((L_{ij}/m_i))\le1$ says that the combined feedback is no more than critical. This can be checked from technologies, adjustment penalties, or preference primitives, provided differentiability, optimality, and the uniform dual-norm estimates really hold. Existence at equality is distinct from stability: optimizing behavior can be mutually consistent even when iterative adjustment keeps oscillating.

\paragraph{Why function-valued choices matter.}
Many economic actions are not scalars or finite vectors. A firm chooses a production trajectory, an investor a state-contingent allocation, a household a consumption or savings path, a planner a policy function, and an informed player a contingent strategy measurable with respect to available information. Such decisions naturally belong to function spaces, and feasible profiles may be bounded and convex without being compact in the norm topology. Reflexive function spaces supply weak compactness of bounded closed convex feasible sets; the distinct gain from (FP) is that the relevant response map need not be \emph{weakly continuous}. Nonexpansiveness in the actual model norm is enough. A separate measurability argument is still necessary for Bayesian or continuum-agent strategies, and choices must respect information constraints. Standard Hilbert and uniformly convex $L^p$ spaces were already covered by older nonexpansive fixed-point results. The claimed increment in generality arises for arbitrary reflexive norms not previously covered by geometric conditions, not simply from the appearance of an infinite-dimensional strategy vector.

\paragraph{Countably many interacting decision makers.}
In \cref{thm:countable}, each decision maker may herself choose an infinite-dimensional plan; the entire profile belongs to a weighted $\ell^p$-sum with $1<p<\infty$. The summability condition is economically more than a technicality: it bounds the aggregate size of feasible plans and makes the complete profile space reflexive. Likewise, the norm bound on the infinite interaction matrix requires collective strategic influence to be non-amplifying in that profile norm. No universal existence statement follows for unrestricted Cartesian products or for a profile space modeled by a nonreflexive $\ell^\infty$ norm. If individual strategy sets are finite-dimensional and tail bounds imply norm compactness, a traditional compact-domain argument might already suffice; one must compare assumptions before declaring a particular application novel.

\paragraph{Aggregate expectations and decentralized consistency.}
In aggregative economies a macroeconomic forecast $m$ or aggregate path affects individual optimal choices; those choices in turn generate a new aggregate $A(R(m))$. The consistency equation $m=A(R(m))$ has the same form as a fixed point. An important point is that nonexpansiveness can be imposed on the \emph{composite} aggregation--response operator, even if neither component is individually nonexpansive. \Cref{cor:aggregate} provides an abstract existence criterion, whereas \cref{thm:nonatomic} supplies a fully specified Bochner-space version for a continuum of negligible agents. In nonatomic games, a player's optimization treats the aggregate as given because a unilateral deviation affects a null set. For finitely many aggregative players, by contrast, own effects on the aggregate must be internalized in the definition of response. A fixed point of a forecast-update equation is not automatically a competitive equilibrium: individual optimization, information, feasibility, budget constraints, and market clearing must all be encoded correctly.

\paragraph{Shared constraints and institutional feasibility.}
The generalized Nash formulation in \cref{thm:generalized} captures settings in which one's admissible decisions depend on other players' actions: shared capacity, an emissions cap, borrowing arrangements, or access to a common resource. At a fixed point the chosen action must both satisfy the endogenous feasible set and maximize the relevant payoff over that set. The theorem avoids imposing a potential structure, but does not relieve the analyst of constructing an everywhere feasible, optimal, norm-nonexpansive selection. In particular, projection operators that are nonexpansive in Hilbert spaces need not remain nonexpansive in arbitrary Banach norms. Invoking projected-gradient or projected-best-response algorithms thus requires a new estimate in the specific geometry used.

\paragraph{Tracking games and the meaning of new generality.}
The target-matching construction in \cref{thm:equivalence} is instructive because it gives a \emph{logical equivalence}, not just a loose analogy. One agent tracks the image $Ty$ of the other's choice while the other tracks the first agent. Mutual optimality forces $x=y=Tx$. Consequently, universal existence in this family is exactly as strong as the nonexpansive fixed-point property. The payoffs are norm-continuous and concave in own actions, yet weak continuity of the induced response may fail. This transfers the full stated geometric generality of (FP) to game theory. It does \emph{not}, by itself, identify a named economic model whose previously open equilibrium problem has been resolved; that stronger historical claim requires independent examination of the model and literature.

\paragraph{Selection, multiplicity, and geometry of the equilibrium set.}
When the additional Bruck-type retraction conclusion (BR) applies, \cref{thm:retraction} supplies a norm-nonexpansive map $P$ from arbitrary reference profiles into the selected equilibrium set, with $P(e)=e$ at equilibria. This yields a mathematically stable \emph{within-game} selection: nearby reference profiles can be assigned nearby equilibria. It does not imply that the chosen equilibrium is nearest to the reference point, welfare-maximizing, unique, computable, or the outcome of a specified learning rule. The retraction also forces the selected equilibrium set to be contractible, hence path-connected. With unique individual best responses, this is a restriction on the \emph{entire} Nash set; with a nonunique correspondence and only a chosen selection $B$, other equilibria outside $\Fix(B)$ may have different topology. Pairwise commuting scenario response operators can have a common fixed point (\cref{cor:commuting}), but commutativity is a substantive compatibility restriction rather than a generic property of policy shocks or unilateral updates.

\paragraph{Existence, dynamics, and approximate equilibrium are different conclusions.}
The isometric best-response cycle in \cref{prop:cycles} shows that an equilibrium may be unique while naive simultaneous best-response learning never converges. Nonexpansiveness means that differences do not expand; it does not mean that errors shrink. On the positive side, anchored contractions yield unique regularized profiles with fixed-point residual at most $\epsilon\operatorname{diam}(C)$, and under uniform own-action payoff Lipschitz estimates, \cref{prop:approx} converts those residuals into explicit regret bounds. This provides an approximation device, not a convergence guarantee for unregularized behavior or a complete algorithmic complexity result. Weak accumulation of approximate fixed points alone is insufficient where weak continuity is absent.

\paragraph{The critical threshold and scope of economic claims.}
The shift-game counterexample of \cref{prop:sharp} makes the threshold precise: a response Lipschitz constant arbitrarily close to but strictly larger than one cannot guarantee equilibrium even with Hilbert-space strategy sets and continuous, strongly own-action-concave payoffs. Hence the nonexpansive hypothesis cannot simply be weakened to ``almost nonexpansive'' without other assumptions. Nor does (FP) imply an equilibrium theorem for arbitrary best-response correspondences, unrestricted mixed-strategy distributions, or nonreflexive spaces such as the standard $L^1$, $L^\infty$, and $C([0,1])$. Restricting mixed strategies to bounded $L^p$ densities can produce an equilibrium only for that restricted deviation set unless additional arguments extend optimality. Finally, Nash existence is silent on Pareto efficiency, market power, fairness, or welfare: stability of a fixed-point equation is not an economic welfare theorem.

\paragraph{Research significance.}
The most promising substantive application is a model with economically interpretable function-valued strategies and a provably nonexpansive optimal-response operator in a \emph{specific original reflexive norm} for which standard weak-continuity and earlier geometric fixed-point arguments do not apply. There the hypotheses translate a simple idea---strategic feedback does not amplify deviations---into an exact existence conclusion at a critical boundary. The remaining research work is to verify the bound from primitives, identify whether the geometry actually lies outside classical existence classes, and separate existence from comparative statics, learning, welfare, and computation. These interpretations are deductions and modeling guidance conditional on (FP) and, where explicitly invoked, (BR), not independent verification of the source theorem.

\section{Research directions and conclusion}
The strongest general consequence is the transfer of \cref{thm:FP} to all games with a globally nonexpansive selected best-response self-map on a bounded convex reflexive domain. The tracking-game equivalence, \cref{thm:equivalence}, shows this is not merely a restatement of finite-dimensional Nash existence: its unrestricted form is logically equivalent to the foundational fixed-point assertion. The spectral-radius and curvature theorems provide primitive, testable sufficient conditions. Infinite-dimensional and continuum-agent variants follow when their full profile spaces are reflexive and their response bounds are genuinely valid. Under the additional Bruck-type source corollary, selected equilibrium sets admit nonexpansive retractions and are contractible. None of these deductions by itself certifies that a specific previously unresolved named economic model has been solved. Such a claim would require verifying a model's best-response structure and checking prior literature separately.

\begingroup\sloppy
\begin{thebibliography}{99}
\bibitem{OAIFP} OpenAI (2026). \emph{Fixed Points of Nonexpansive Maps in Reflexive Banach Spaces}, 24 September 2026. \href{https://github.com/openai/math/blob/adc7f1241b42e322a6451854ab7e4b4c146bf78a/preprints/Fixed-Points-of-Nonexpansive-Maps-in-Reflexive-Banach-Spaces-September-24-2026/paper.pdf}{manuscript PDF}.
\bibitem{Tracker} Erdos Problems LLM Hunter, problem LOCAL-30007184, ``Nonexpansive fixed points in reflexive Banach spaces.'' \href{https://mehmetmars7.github.io/Erdosproblems-llm-hunter/problem.html?type=open_problems&id=30007184}{problem page}.
\bibitem{LeanMain} OpenAI (2026). \texttt{lean/OAI/Analysis/Nonexpansive/Main.lean}. \href{https://github.com/openai/math/blob/adc7f1241b42e322a6451854ab7e4b4c146bf78a/lean/OAI/Analysis/Nonexpansive/Main.lean}{Lean implementation}.
\bibitem{LeanScope} OpenAI (2026). Formalization scope, family 328. \href{https://github.com/openai/math/blob/adc7f1241b42e322a6451854ab7e4b4c146bf78a/lean/docs/328.md}{formalization scope}.
\bibitem{LeanComparator} OpenAI (2026). \texttt{lean/ComparatorChallenges/ReflexiveFixedPoints.lean} and its companion \texttt{.json} configuration. \href{https://github.com/openai/math/blob/adc7f1241b42e322a6451854ab7e4b4c146bf78a/lean/ComparatorChallenges/ReflexiveFixedPoints.lean}{Comparator statement}.
\bibitem{Glicksberg1952} I. L. Glicksberg (1952). A further generalization of the Kakutani fixed point theorem, with application to Nash equilibrium points. \emph{Proceedings of the American Mathematical Society} 3, 170--174.
\bibitem{Browder1965} F. E. Browder (1965). Nonexpansive nonlinear operators in a Banach space. \emph{Proceedings of the National Academy of Sciences USA} 54, 1041--1044.
\bibitem{Gohde1965} D. G\"ohde (1965). Zum Prinzip der kontraktiven Abbildung. \emph{Mathematische Nachrichten} 30, 251--258.
\bibitem{Kirk1965} W. A. Kirk (1965). A fixed point theorem for mappings which do not increase distances. \emph{American Mathematical Monthly} 72, 1004--1006.
\bibitem{Bruck1974} R. E. Bruck, Jr. (1974). A common fixed point theorem for a commuting family of nonexpansive mappings. \emph{Pacific Journal of Mathematics} 53, 59--71. \url{https://msp.org/pjm/1974/53-1/pjm-v53-n1-p05-p.pdf}.
\end{thebibliography}
\endgroup
\end{document}
